\section{Dinâmica Relativística e Escala Gamma}
\label{sec:relativistica}

O segundo fundamento conceitual vem da formulação de DRM com dinâmica não
linear dirigida por escala relativística \cite{muniz_relativistic_drm}. Nessa
formulação, um mapa iterativo
\begin{equation}
  x_{n+1}=F(x_n; A(\omega,r), D(x_n), g_{x_n})
\end{equation}
é controlado por uma amplitude
\begin{equation}
  A(\omega,r) = \frac{4}{\pi}\omega r \gamma(\omega r),
\end{equation}
onde $\gamma$ é o fator de Lorentz e $v=\omega r$ é uma velocidade tangencial.
À medida que $v$ se aproxima de $c$, $\gamma(v)$ diverge:
\begin{equation}
  \gamma(v)=\frac{1}{\sqrt{1-v^2/c^2}}.
\end{equation}

No paper de dinâmica relativística, essa divergência é associada a regimes de
bifurcação, multiatratores e, sob hipóteses explícitas, convergência para
atratores toroidais. O DRM Transformer não implementa literalmente o mapa
logístico direcional do paper. Em vez disso, usa a intuição da escala
relativística como modulação de resolução geométrica na atenção.

Na implementação, cada token projetado no manifold tem distância a âncoras
semânticas. Essa distância gera um fator gamma suavizado, limitado e aplicado à
distância geodésica:
\begin{equation}
  \tilde{d}_{ij}^2 = d_{G,ij}^2
    \left(1+\alpha \log(1+\gamma_i-1)\right)^2.
\end{equation}
O parâmetro $\alpha$ controla a força do gamma-scaling. O uso do logaritmo e de
clamps evita singularidades numéricas e preserva estabilidade durante o treino.
Assim, o fator relativístico funciona como um mecanismo adaptativo de escala,
sem reivindicar nova física fundamental.
